\section*{Chapitre \ref{ch:g12:reaction-rates} --- La vitesse des réactions}

\begin{solution}{exo:g12:reaction-rates:1}
Température~; surface de contact~; concentration~; température.
\end{solution}

\begin{solution}{exo:g12:reaction-rates:2}
\qty{13.9}{min} (en bleu) et \qty{6.9}{min} (en rouge). Au bout de deux
\omterm{def:g12:reaction-rates:half-life}{temps de demi-réaction},
$10.0 / 4 = \qty{2.5}{mmol/L}$.
\end{solution}

\begin{solution}{exo:g12:reaction-rates:3}
Pente $(5.0 - 1.0)/10 = \qty{0.40}{mmol.L^{-1}.min^{-1}}$~: c'est la
\omterm{def:g12:reaction-rates:rate}{vitesse d'apparition} à \qty{5}{min}.
\end{solution}

\begin{solution}{exo:g12:reaction-rates:4}
Le refroidissement et la \omterm{def:g10:concentration-and-dilution:dilution}{dilution} rendent la réaction si lente que la
composition ne change plus pendant le \omterm{def:g11:titration:titration}{titrage}. La composition mesurée
est celle qu'avait le \omterm{def:g6:pure-substances-and-mixtures:mixture}{mélange} quand la réaction a été arrêtée, au moment
où on l'a versé.
\end{solution}

\begin{solution}{exo:g12:reaction-rates:5}
\qty{30}{min} représentent 3 \omterm{def:g12:reaction-rates:half-life}{temps de demi-réaction}~: il reste $1/8$.
\qty{1}{h} en représente 6~:
$1/64$, soit environ \qty{1.6}{\percent}.
\end{solution}

\begin{solution}{exo:g12:reaction-rates:6}
$\ln[\ce{A}]$~: 2.079, 1.733, 1.386, 1.040, 0.693. Il diminue de 0.346
tous les
\qty{5}{min}~: c'est une droite, la réaction est d'ordre 1. $k = 0.346/5 =
\qty{0.0693}{min^{-1}}$ et $t_{1/2} = \ln 2 / k = \qty{10.0}{min}$ (ce
qui se voit directement~: 8.00, 4.00, 2.00 tous les \qty{10}{min}).
\end{solution}

\begin{solution}{exo:g12:reaction-rates:7}
$k = 0.693 / 25 = \qty{0.0277}{s^{-1}}$.
\end{solution}

\begin{solution}{exo:g12:reaction-rates:8}
$0.100\,e^{-0.20} = \qty{0.0819}{mol/L}$; $0.100\,e^{-0.60} =
\qty{0.0549}{mol/L}$; $0.100\,e^{-2.0} = \qty{0.0135}{mol/L}$.
\end{solution}

\begin{solution}{exo:g12:reaction-rates:9}
À \qty{13.9}{min}, $[\ce{A}] = \qty{5.0}{mmol/L}$~; la pente de la
tangente est proche de $-0.25$, et $-k[\ce{A}] = -0.050 \times 5.0 =
\qty{-0.25}{mmol.L^{-1}.min^{-1}}$~: la moitié de la valeur initiale.
\end{solution}

\begin{solution}{exo:g12:reaction-rates:10}
Le diiode est la seule espèce colorée~: d'après la loi de Beer--Lambert,
$A = \varepsilon \ell [\ce{I2}]$. La \omterm{def:g12:reaction-rates:rate}{vitesse d'apparition} est la pente
de la tangente à la courbe $[\ce{I2}](t)$, c'est-à-dire la pente de la
tangente à $A(t)$ divisée par $\varepsilon \ell$.
\end{solution}

\begin{solution}{exo:g12:reaction-rates:11}
La vitesse initiale $k[\ce{A}]_0$ est deux fois plus grande pour la
seconde expérience~; les \omterm{def:g12:reaction-rates:half-life}{temps de demi-réaction} sont égaux,
$\ln 2 / k$, indépendants de $[\ce{A}]_0$.
\end{solution}

\begin{solution}{exo:g12:reaction-rates:12}
$t = \ln(1/f)/k$. En \omterm{def:g12:reaction-rates:half-life}{temps de demi-réaction}~: $\ln(1/f)/\ln 2$. Pour
\qty{1}{\percent},
$\ln 100 / \ln 2 = 6.6$~; pour \qty{0.1}{\percent}, $\ln 1000 / \ln 2 =
10.0$.
\end{solution}

\begin{solution}{exo:g12:reaction-rates:13}
\qty{90}{\percent} consommés, $f = 0.10$~: $t = \ln 10 / k$, soit
\qty{46}{min} à \qty{20}{\celsius} et \qty{23}{min} à
\qty{30}{\celsius}. La \omterm{def:g12:reaction-rates:first-order}{constante de vitesse} double sur
\qty{10}{\celsius}~: la règle est vérifiée ici.
\end{solution}

\begin{solution}{exo:g12:reaction-rates:14}
La température descend vers \qty{0}{\celsius}, et les concentrations
sont divisées par 10. Avec une vitesse proportionnelle au produit de
deux concentrations, la seule \omterm{def:g10:concentration-and-dilution:dilution}{dilution} la divise par
$10 \times 10 = 100$.
\end{solution}

\begin{solution}{exo:g12:reaction-rates:15}
$v(t) = -\dfrac{d}{dt}\big([\ce{A}]_0 e^{-kt}\big) = k[\ce{A}]_0 e^{-kt}$.
Vitesse initiale $k[\ce{A}]_0$~; la vitesse est divisée par deux quand
$e^{-kt} = 1/2$, c'est-à-dire au bout d'un \omterm{def:g12:reaction-rates:half-life}{temps de demi-réaction}.
\end{solution}

\begin{solution}{pb:g12:reaction-rates:1}
\textbf{1.} D'après la loi de Beer--Lambert, le colorant étant la seule
espèce qui absorbe à cette longueur d'onde~; la solution doit être assez
diluée ($A$ inférieure à 1 environ).

\textbf{2.} Il n'absorbe pas la lumière visible~: il n'ajoute rien à
$A$.

\textbf{3.} Pour que leur concentration varie à peine~: la vitesse ne
dépend alors que du colorant.

\textbf{4.} $A = 0$ (le blanc).

\textbf{5.} $A$ passe de 0.800 à 0.402 en \qty{6}{min}~: environ
\qty{6}{min}.

\textbf{6.} De 0.402 à \qty{6}{min} à 0.199 à \qty{12}{min}~: de
nouveau divisée par deux en \qty{6}{min}.

\textbf{7.} $-0.223$, $-0.448$, $-0.691$, $-0.911$, $-1.155$, $-1.366$,
$-1.614$, $-1.945$, $-2.551$, $-3.079$.

\textbf{8.} Les points sont alignés~: la réaction est d'ordre 1 par
rapport au colorant.

\textbf{9.} $(-2.551 + 0.223)/20 = \qty{-0.116}{min^{-1}}$.

\textbf{10.} $k = \qty{0.116}{min^{-1}}$.

\textbf{11.} $0.693 / 0.116 = \qty{6.0}{min}$, comme on l'a lu dans le
tableau.

\textbf{12.} $k A_0 = 0.116 \times 0.800 = \qty{0.093}{min^{-1}}$ en
unités d'\omterm{def:g11:absorbance:absorbance}{absorbance} par minute.

\textbf{13.} Le même, \qty{6.0}{min}~: le \omterm{def:g12:reaction-rates:half-life}{temps de demi-réaction} d'une
\omterm{def:g12:reaction-rates:first-order}{réaction d'ordre 1} ne dépend pas de la valeur de départ.

\textbf{14.} $0.800\, e^{-0.116 \times 30} = 0.025$.

\textbf{15.} $t = \ln 100 / k = 4.61 / 0.116 = \qty{40}{min}$.

\textbf{16.} $40 / 6.0 = 6.6$ \omterm{def:g12:reaction-rates:half-life}{temps de demi-réaction}.

\textbf{17.} $t = \ln(0.800/0.010)/k = 4.38 / 0.116 = \qty{38}{min}$.

\textbf{18.} Toujours 6.6 \omterm{def:g12:reaction-rates:half-life}{temps de demi-réaction}, soit maintenant
$6.6 \times 3.0 = \qty{20}{min}$.

\textbf{19.} Environ 6.6 \omterm{def:g12:reaction-rates:half-life}{temps de demi-réaction}, soit environ
\qty{40}{min} ici.
\end{solution}
